\chapter{Background and Related Work}
\section{Classical Approaches to Inventory Control}
\label{sec:classical-inventory-control}

The study of inventory control under uncertainty begins with the single-period newsvendor
model, which balances the cost of holding unsold stock against the cost of unmet demand and
places the optimal order quantity at the critical fractile of the demand distribution
[CITE: newsvendor model and critical fractile, e.g. Zipkin 2000 or Silver, Pyke and Thomas 2017].
Across multiple periods this trade-off becomes the dynamic stochastic inventory problem
[CITE: Arrow, Harris and Marschak 1951 as the foundational dynamic stochastic inventory work],
for which the structure of the optimal replenishment policy is well understood. When each
replenishment incurs a fixed ordering cost, the optimal policy takes the $(s,S)$ form and places
an order to raise the inventory position to $S$ whenever it falls to or below $s$
[CITE: Scarf 1960 optimality of (s,S) policies]. When no fixed ordering cost is present, this
reduces to a base-stock or order-up-to policy that raises the inventory position to a single
target level each period [CITE: base-stock policy optimality, e.g. Veinott 1966 and Zipkin 2000].
The deterministic special case in which demand is constant is the classical economic order
quantity model [CITE: Harris 1913 EOQ and Erlenkotter 1990 historical account]. The network
studied in this thesis carries no fixed ordering cost, so the base-stock policy is the relevant
optimal class and underlies the reference policies it evaluates.

A second strand of the classical literature concerns how the parameters of a chosen policy are
computed. When demand follows a tractable distribution and unmet demand is backordered,
base-stock levels admit closed-form or readily computable expressions
[CITE: analytical base-stock levels, Zipkin 2000]. This tractability is lost under the lost-sales
assumption adopted in this thesis, where unmet demand disappears rather than carrying over and the
optimal policy has no simple closed form [CITE: Zipkin 2008 old and new methods for lost-sales
systems], a setting whose reliance on approximations and numerical procedures is documented in
dedicated reviews [CITE: Bijvank and Vis 2011 lost-sales review]. Two practical responses recur.
The first calibrates policy parameters analytically from estimated demand statistics, as in
newsvendor-style safety-stock formulas. The second treats the simulator as a black box and searches
for the parameters that minimise simulated cost, an approach known as simulation optimisation
[CITE: simulation optimisation for inventory base-stock levels]. Forecast-driven and adaptive
heuristics that update their parameters from observed demand occupy the middle ground between the
two [CITE: adaptive or forecast-based replenishment heuristics, Silver, Pyke and Thomas 2017]. The
reference policies used in this thesis instantiate this spectrum from analytically derived to
simulation-optimised base-stock levels.

Moving from a single location to a network introduces coordination across stocking points. The
multi-echelon literature, beginning with the analysis of serial systems, establishes that echelon
base-stock policies are optimal under suitable conditions and provides the foundation for inventory
control across tiered supply chains [CITE: Clark and Scarf 1960 multi-echelon serial systems].
Related work extends these ideas to repairable-item and service-parts networks
[CITE: Sherbrooke 1968 METRIC]. The setting closest to the present work is the single-echelon,
multi-location distribution system in which warehouses serve a set of demand regions, for which the
optimal allocation policy is generally intractable and is addressed through heuristics
[CITE: Eppen and Schrage 1981 and Federgruen and Zipkin 1984 one-warehouse multi-retailer]. Within
such networks, lateral transshipment allows stock to move between locations of the same echelon to
cover demand that cannot be met locally [CITE: Paterson et al. 2011 lateral transshipment review].
The environment in this thesis realises a reactive form of this mechanism implicitly through its
greedy fulfilment of demand from the cheapest stocked warehouse. Comprehensive surveys catalogue
the assumptions and methods of this body of work and note that most models assume a single echelon
or a strictly tiered structure [CITE: de Kok et al. 2018 typology and review of stochastic
multi-echelon inventory models].

Two features limit the reach of these classical methods in the setting this thesis studies. First,
their tractability rests on structural and distributional assumptions that rarely hold in realistic
networks, and the dynamic programmes that characterise optimal policies suffer from the curse of
dimensionality as the number of locations and products grows [CITE: curse of dimensionality in
inventory networks, e.g. de Kok et al. 2018]. Optimal policies for general distribution systems are
therefore unknown, and practice relies on heuristics whose quality is difficult to certify. Second,
the analytically convenient cases are typically symmetric across locations and assume centralised
information, whereas the networks of interest here are heterogeneous and operate under decentralised
local information. These limitations motivate data-driven and learning-based methods, which
approximate replenishment policies directly from interaction with the system rather than from
closed-form analysis [CITE: Boute, Gijsbrechts, van Jaarsveld and Vanvuchelen 2022 DRL for
inventory control roadmap]. The remainder of this chapter reviews reinforcement learning as such an
approach.</parameter>

------------------------------------------------------------------------------------------------

Research on stochastic inventory control originates with the single-period newsvendor model, in
which a single order is placed before demand is observed and the order quantity balances the
expected cost of overage against that of underage. The cost-minimising quantity equates the
marginal costs of the two errors and is given by the critical fractile of the demand distribution
[CITE: newsvendor model and critical fractile, e.g. Zipkin 2000]. Arrow, Harris and Marschak (1951)
[CITE: Arrow, Harris and Marschak 1951, Optimal Inventory Policy] generalise this static decision
into a multi-period dynamic programme and give the first integrated account of inventory decisions
under sequential demand uncertainty. Scarf (1960) [CITE: Scarf 1960, optimality of (s,S) policies]
characterises the optimal policy for the periodic-review problem with a fixed ordering cost. He
introduces the notion of $K$-convexity and uses it to prove that an $(s,S)$ policy is optimal,
under which an order raising the inventory position to $S$ is placed whenever the position drops to
or below $s$. Veinott (1966) [CITE: Veinott 1966, conditions for (s,S) optimality] later weakens
the conditions required for this result. In the absence of a fixed ordering cost the reorder point
and the order-up-to level coincide, and the optimum reduces to a base-stock policy with a single
target level [CITE: base-stock optimality, Veinott 1966 and Zipkin 2000]. The deterministic
antecedent of this work is the economic order quantity of Harris (1913) [CITE: Harris 1913 EOQ and
Erlenkotter 1990 historical account], which trades a fixed ordering cost against a linear holding
cost under constant demand and remains the textbook entry point to the field.

A parallel question is how the parameters of these policies are computed once their structure is
fixed. In the backorder case, where unmet demand carries to later periods, base-stock levels follow
from tractable newsvendor-type expressions [CITE: analytical base-stock computation, Zipkin 2000].
The lost-sales case, where unmet demand is simply lost, is markedly harder. Zipkin (2008) [CITE:
Zipkin 2008, lost-sales inventory systems] re-examines it and shows that the optimal policy is no
longer of base-stock form and depends on the full vector of outstanding orders. Bijvank and Vis
(2011) [CITE: Bijvank and Vis 2011, lost-sales review] survey the resulting literature and find that
analysis proceeds almost entirely through restricted policy classes and approximations rather than
optimal rules. Where analytical levels are unavailable, a practical alternative evaluates a
candidate policy by simulation and searches the parameter space for the configuration of lowest
cost, an approach studied under the name of simulation optimisation [CITE: simulation optimisation
for inventory systems]. Adaptive heuristics occupy the ground between these extremes, re-estimating
demand from a moving window and updating their order-up-to levels online [CITE: adaptive or
forecast-based replenishment, Silver, Pyke and Thomas 2017].

Extending these results from one stocking point to a network is the concern of multi-echelon
inventory theory. Clark and Scarf (1960) [CITE: Clark and Scarf 1960, multi-echelon serial systems]
analyse a serial system and show that it decomposes into a sequence of single-stage problems linked
by induced penalty functions, establishing that an echelon base-stock policy, defined on the stock
of a stage together with all stock downstream of it, is optimal. Sherbrooke (1968) [CITE:
Sherbrooke 1968, METRIC] develops the METRIC model for repairable-item networks at the RAND
Corporation and sets stock levels across a two-echelon depot and base structure by approximating
the expected backorders at each base under compound Poisson demand. For distribution systems, in
which a depot supplies several retailers facing independent demand, the optimal policy is generally
intractable. Eppen and Schrage (1981) [CITE: Eppen and Schrage 1981, one-warehouse multi-retailer]
obtain closed-form results under a balance assumption that permits negative allocations, and
Federgruen and Zipkin (1984) [CITE: Federgruen and Zipkin 1984, multilocation approximations]
construct approximations for the general multilocation problem. A closely related mechanism is
lateral transshipment, the movement of stock between locations of the same echelon. Paterson et al.
(2011) [CITE: Paterson et al. 2011, lateral transshipment review] survey this work and separate
proactive transshipments, carried out at predetermined times, from reactive transshipments that
respond to stockouts. The assumptions and solution methods spanning these models are catalogued by
de Kok et al. (2018) [CITE: de Kok et al. 2018, typology and review], whose typology shows that
most tractable formulations rest on a single echelon or a strictly serial structure.

Two limitations run through this body of work. First, its tractable results depend on strong
structural assumptions, and the exact dynamic programmes behind them become computationally
prohibitive as the number of locations, products and lead-time stages grows, a difficulty long
known as the curse of dimensionality [CITE: curse of dimensionality in inventory, e.g. de Kok et
al. 2018]. For general distribution networks the optimal policy is unknown, and the field proceeds
through heuristics whose optimality gap is seldom quantified. Second, the analytically convenient
models are largely symmetric across locations and assume a central planner with full network
information, whereas real distribution networks are heterogeneous across sites and are operated from
local information. These gaps have motivated a shift toward data-driven methods that learn
replenishment policies from interaction rather than deriving them in closed form, a direction
surveyed for inventory control by Boute et al. (2022) [CITE: Boute et al. 2022, DRL for inventory
control roadmap]. The remainder of this chapter reviews the reinforcement learning methods on which
that approach rests.


\section{Reinforcement Learning}
\label{sec:reinforcement-learning}
\subsection{Single-Agent Reinforcement Learning}
\label{sec:single-agent-rl}

\begin{itemize}
    \item Single Agent RL general 
    \item Markov Decision Process
    \item Overview of Approaches
    \item Policy Gradient Approaches and PPO
\end{itemize}


\subsection{Multi-Agent Reinforcement Learning}

\begin{itemize}
    \item General MARL concepts
    \item Types of Games and Dec-POMDP
    \item Challenges (non-stationarity, credit assignment, ...)
    \item Overview of Approaches (Independent Learners, Joint Learners, cdte, parameter sharing, mappo, ippo, evtl. also mention value decomposition, ...)
    \item Policy Gradient Approaches and PPO
\end{itemize}

\section{Heterogeneity in Multi-Agent Reinforcement Learning}

\section{2.4  Research Gap and Research Questions}

\paragraph{\textbf{G1. Studied inventory control problems are in most cases symmetric.}}

\paragraph{\textbf{G2. No framework links heterogeneity types to failure modes.}}

\paragraph{\textbf{G3. Existing remedies are untargeted.}}

\paragraph{\textbf{RQ1}} 
\textit{While cooperative multi-agent reinforcement learning with parameter sharing achieves strong results on symmetric multi-warehouse environments, real-world supply chain settings are inherently heterogeneous. Under such heterogeneity, does parameter-shared multi-agent PPO degrade relative to a symmetric baseline of matched network throughput and if so, what failure modes explain the gap? Can environment-specific hyperparameter tuning recover this lost performance?}

\paragraph{\textbf{RQ2}}  
\textit{Parameter sharing imposes a single inductive bias across all agents, which becomes a liability when agents operate in heterogeneous conditions. Which architectural modifications relax this constraint sufficiently to recover performance and are the same modifications effective across all heterogeneity types or does each heterogeneity type demand a distinct architectural response?}
